
% Overleaf 编译建议：
% Menu -> Settings -> Compiler -> XeLaTeX
% 本文使用 Overleaf/TeX Live 自带的 Fandol 中文字体，无需上传任何字体文件。
\documentclass[UTF8,12pt,a4paper,fontset=fandol]{ctexart}

\usepackage{amsmath,amssymb,bm}
\usepackage{geometry}
\usepackage{booktabs}
\usepackage{hyperref}
\hypersetup{
    unicode=true,
    colorlinks=true,
    linkcolor=black,
    urlcolor=black,
    citecolor=black
}
\usepackage{enumitem}

\geometry{left=2.4cm,right=2.4cm,top=2.5cm,bottom=2.5cm}
\setlength{\parindent}{2em}
\setlength{\parskip}{0.4em}

\title{有限脉冲一阶 Bragg 原子干涉仪传递函数计算\\
\large 平面波、无光学谐振腔情形}
\author{}
\date{}

\begin{document}
\maketitle

\section{问题设定}

本文讨论一个三脉冲 Mach--Zehnder 型 Bragg 原子干涉仪，采用两束反向传播平面波，不考虑光学谐振腔，只考虑一阶 Bragg 衍射，并考虑有限脉冲宽度。

脉冲序列为
\[
\frac{\pi}{2}-\pi-\frac{\pi}{2}.
\]

记 Bragg 有效 Rabi 频率为
\[
f_{\rm R},
\]
若其单位为 Hz，则对应的角 Rabi 频率为
\[
\boxed{\Omega_{\rm R}=2\pi f_{\rm R}}.
\]

本文的目标是计算光学相位扰动
\[
\delta\varphi(t)
\]
到原子干涉仪输出相位
\[
\delta\Phi_{\rm AI}(t)
\]
之间的频域传递函数。

MIGA 文章采用 sensitivity function 形式写原子干涉仪响应：
\[
\Delta\phi_{\rm AT}(t)
=
n\,\Delta\varphi(t)*\frac{ds(t)}{dt},
\]
其中 $s(t)$ 是三脉冲原子干涉仪的 sensitivity function，$*$ 表示卷积。对于本文的一阶 Bragg，取
\[
n=1.
\]

因此
\[
\boxed{
\Delta\Phi_{\rm AI}(t)
=
\Delta\varphi(t)*\frac{ds(t)}{dt}
}.
\]

需要说明：MIGA 文章本身没有给出本文所用有限脉冲 $s(t)$ 的显式分段表达式，而是引用 sensitivity-function 文献。下文是在该 formalism 下，将一阶 Bragg 近似为共振两能级 Rabi 演化得到的有限脉冲结果。

\section{脉冲宽度与时间轴}

对共振方波脉冲，脉冲面积满足
\[
\theta=\Omega_{\rm R}\tau.
\]

对于 $\pi/2$ 脉冲：
\[
\Omega_{\rm R}\tau=\frac{\pi}{2},
\]
所以
\[
\boxed{
\tau=\frac{\pi}{2\Omega_{\rm R}}
=\frac{1}{4f_{\rm R}}
}.
\]

因此 $\pi$ 脉冲持续时间为
\[
\boxed{
\tau_\pi=2\tau=\frac{1}{2f_{\rm R}}
}.
\]

本文采用如下时间定义：

\begin{itemize}[leftmargin=2em]
\item 第一束 $\pi/2$ 脉冲：
\[
0\le t<\tau;
\]
\item 第一段自由演化：
\[
\tau\le t<T+\tau;
\]
\item 第二束 $\pi$ 脉冲：
\[
T+\tau\le t<T+3\tau;
\]
\item 第二段自由演化：
\[
T+3\tau\le t<2T+3\tau;
\]
\item 第三束 $\pi/2$ 脉冲：
\[
2T+3\tau\le t<2T+4\tau.
\]
\end{itemize}

在这个定义中，$T$ 是两段自由演化时间，因此整个脉冲序列持续时间为
\[
\boxed{
T_{\rm total}=2T+4\tau
}.
\]

如果实验中把 $T$ 定义为相邻脉冲中心之间的时间，则只需要重新平移各脉冲的起止时刻，计算方法不变。

\section{有限脉冲 sensitivity function}

对于理想瞬时脉冲，$s(t)$ 是分段常数；对于有限脉冲，原子在脉冲内部发生连续 Rabi rotation，因此 $s(t)$ 在脉冲期间连续变化。

在本文的共振两能级近似下，可写为
\[
\boxed{
s(t)=
\begin{cases}
0,
&t<0,
\\[2mm]
\sin(\Omega_{\rm R}t),
&0\le t<\tau,
\\[2mm]
1,
&\tau\le t<T+\tau,
\\[2mm]
\cos\!\left[\Omega_{\rm R}(t-T-\tau)\right],
&T+\tau\le t<T+3\tau,
\\[2mm]
-1,
&T+3\tau\le t<2T+3\tau,
\\[2mm]
-\cos\!\left[\Omega_{\rm R}(t-2T-3\tau)\right],
&2T+3\tau\le t<2T+4\tau,
\\[2mm]
0,
&t\ge 2T+4\tau.
\end{cases}
}
\]

该表达式满足正确的边界条件：

\[
s(0)=0,
\qquad
s(\tau)=1,
\]
第二个 $\pi$ 脉冲使
\[
1\rightarrow -1,
\]
最后一个 $\pi/2$ 脉冲使
\[
-1\rightarrow 0.
\]

\section{由 sensitivity function 定义传递函数}

定义 $s(t)$ 的 Fourier transform 为
\[
\boxed{
G(\omega)
=
\int_{-\infty}^{+\infty}
s(t)e^{-i\omega t}\,dt
}.
\]

由于
\[
\Delta\Phi_{\rm AI}(t)
=
\Delta\varphi(t)*\frac{ds(t)}{dt},
\]
对其做 Fourier transform：
\[
\Delta\Phi_{\rm AI}(\omega)
=
\Delta\varphi(\omega)\,
\mathcal{F}
\left[
\frac{ds}{dt}
\right].
\]

利用
\[
\mathcal{F}
\left[
\frac{ds}{dt}
\right]
=
i\omega G(\omega),
\]
得到
\[
\boxed{
\Delta\Phi_{\rm AI}(\omega)
=
H_\varphi(\omega)
\Delta\varphi(\omega)
}
\]
其中
\[
\boxed{
H_\varphi(\omega)
=
i\omega G(\omega)
}.
\]

因此相位传递函数的幅值为
\[
\boxed{
|H_\varphi(\omega)|
=
|\omega G(\omega)|
}
\]
而功率谱传递因子为
\[
\boxed{
|H_\varphi(\omega)|^2
=
|\omega G(\omega)|^2
}.
\]

\section{直接计算 $\dot s(t)$}

数值上，比起先计算 $G(\omega)$ 再乘 $i\omega$，更方便的方法是直接对
\[
\dot s(t)=\frac{ds}{dt}
\]
做 Fourier transform。

由上面的 $s(t)$ 可得
\[
\boxed{
\dot s(t)=
\begin{cases}
\Omega_{\rm R}\cos(\Omega_{\rm R}t),
&0<t<\tau,
\\[2mm]
0,
&\tau<t<T+\tau,
\\[2mm]
-\Omega_{\rm R}
\sin\!\left[\Omega_{\rm R}(t-T-\tau)\right],
&T+\tau<t<T+3\tau,
\\[2mm]
0,
&T+3\tau<t<2T+3\tau,
\\[2mm]
\Omega_{\rm R}
\sin\!\left[\Omega_{\rm R}(t-2T-3\tau)\right],
&2T+3\tau<t<2T+4\tau,
\\[2mm]
0,
&\text{其他时刻}.
\end{cases}
}
\]

于是可以直接写
\[
\boxed{
H_\varphi(\omega)
=
\int_{-\infty}^{+\infty}
\dot s(t)e^{-i\omega t}\,dt
}.
\]

这就是有限脉冲一阶 Bragg 干涉仪的相位传递函数。

\section{分段 Fourier 积分}

若希望解析计算，则可以写
\[
G(\omega)=G_1+G_2+G_3+G_4+G_5,
\]
其中

\subsection*{第一束 $\pi/2$ 脉冲}
\[
G_1
=
\int_0^\tau
\sin(\Omega_{\rm R}t)e^{-i\omega t}\,dt.
\]

\subsection*{第一段自由演化}
\[
G_2
=
\int_\tau^{T+\tau}
e^{-i\omega t}\,dt
=
\frac{
e^{-i\omega\tau}
-
e^{-i\omega(T+\tau)}
}{
i\omega
}.
\]

\subsection*{第二束 $\pi$ 脉冲}
\[
G_3
=
\int_{T+\tau}^{T+3\tau}
\cos\!\left[\Omega_{\rm R}(t-T-\tau)\right]
e^{-i\omega t}\,dt.
\]

令
\[
u=t-(T+\tau),
\]
则
\[
G_3
=
e^{-i\omega(T+\tau)}
\int_0^{2\tau}
\cos(\Omega_{\rm R}u)e^{-i\omega u}\,du.
\]

\subsection*{第二段自由演化}
\[
G_4
=
-\int_{T+3\tau}^{2T+3\tau}
e^{-i\omega t}\,dt.
\]

\subsection*{第三束 $\pi/2$ 脉冲}
\[
G_5
=
-\int_{2T+3\tau}^{2T+4\tau}
\cos\!\left[\Omega_{\rm R}(t-2T-3\tau)\right]
e^{-i\omega t}\,dt.
\]

最后
\[
\boxed{
H_\varphi(\omega)
=
i\omega
\left(
G_1+G_2+G_3+G_4+G_5
\right)
}.
\]

实际做数据分析或仿真时，更推荐直接数值积分，而不是手工化简上面的解析式。

\section{推荐的数值计算流程}

给定
\[
f_{\rm R},\qquad T,
\]
首先计算
\[
\Omega_{\rm R}=2\pi f_{\rm R},
\qquad
\tau=\frac{1}{4f_{\rm R}}.
\]

然后构造时间数组 $t_j$，根据前述分段公式计算 $s(t_j)$ 或 $\dot s(t_j)$。

对于每一个目标频率
\[
f,
\qquad
\omega=2\pi f,
\]
直接做
\[
\boxed{
H_\varphi(\omega)
=
\int
\dot s(t)e^{-i\omega t}dt
}.
\]

离散形式可写为
\[
H_\varphi(\omega)
\approx
\sum_j
\dot s(t_j)e^{-i\omega t_j}\Delta t.
\]

最终画出
\[
|H_\varphi(f)|
\]
或
\[
|H_\varphi(f)|^2.
\]

\section{两个必须满足的校验条件}

\subsection{零频响应}

因为
\[
H_\varphi(0)
=
\int_{-\infty}^{+\infty}
\dot s(t)\,dt
=
s(+\infty)-s(-\infty),
\]
而序列前后均有
\[
s=0,
\]
所以
\[
\boxed{
H_\varphi(0)=0
}.
\]

这表示一个恒定的整体激光相位偏移不会被 Mach--Zehnder 原子干涉仪检测到。

\subsection{瞬时脉冲极限}

令
\[
f_{\rm R}\rightarrow\infty,
\qquad
\tau\rightarrow0.
\]

此时
\[
s(t)
\rightarrow
\begin{cases}
1,&0<t<T,\\
-1,&T<t<2T,\\
0,&\text{其他}.
\end{cases}
\]

因此
\[
\dot s(t)
=
\delta(t)-2\delta(t-T)+\delta(t-2T).
\]

于是
\[
H_\varphi(\omega)
=
1-2e^{-i\omega T}+e^{-i2\omega T}.
\]

可写为
\[
H_\varphi(\omega)
=
(1-e^{-i\omega T})^2.
\]

所以其幅值为
\[
\boxed{
|H_\varphi(\omega)|
=
4\sin^2\left(\frac{\omega T}{2}\right)
}.
\]

功率谱传递函数为
\[
\boxed{
|H_\varphi(\omega)|^2
=
16\sin^4\left(\frac{\omega T}{2}\right)
}.
\]

因此，当
\[
\omega\ll\Omega_{\rm R}
\]
且脉冲足够短时，有限脉冲数值结果应逐渐逼近该表达式。

\section{一阶 Bragg 的有效波矢}

对于一阶两光子 Bragg 过程，
\[
\boxed{
k_{\rm eff}\simeq 2k
}
\]
其中
\[
k=\frac{2\pi}{\lambda}.
\]

因此
\[
\boxed{
k_{\rm eff}
=
\frac{4\pi}{\lambda}
}.
\]

若一个位移扰动 $x(t)$ 引起的光学相位扰动可写成
\[
\delta\varphi(t)
=
k_{\rm eff}x(t),
\]
则
\[
\Delta\Phi_{\rm AI}(\omega)
=
k_{\rm eff}
H_\varphi(\omega)
x(\omega).
\]

所以位移到原子相位的传递函数为
\[
\boxed{
H_{x\rightarrow\Phi}(\omega)
=
k_{\rm eff}H_\varphi(\omega)
=
k_{\rm eff}i\omega G(\omega)
}.
\]

\section{加速度输入}

若输入为沿 Bragg 轴方向的加速度 $a(t)$，则在频域有
\[
a(\omega)
=
-\omega^2x(\omega),
\]
即
\[
x(\omega)
=
-\frac{a(\omega)}{\omega^2}.
\]

因此
\[
\Delta\Phi_{\rm AI}(\omega)
=
-k_{\rm eff}
\frac{i\omega G(\omega)}{\omega^2}
a(\omega).
\]

于是加速度到原子相位的传递函数为
\[
\boxed{
H_{a\rightarrow\Phi}(\omega)
=
-k_{\rm eff}
\frac{i\omega G(\omega)}{\omega^2}
}.
\]

在低频极限，
\[
\sin\left(\frac{\omega T}{2}\right)
\simeq
\frac{\omega T}{2},
\]
所以瞬时脉冲近似下
\[
|H_{a\rightarrow\Phi}|
\rightarrow
k_{\rm eff}T^2.
\]

这恢复了原子干涉仪的经典结果
\[
\boxed{
\Delta\Phi
=
k_{\rm eff}aT^2
}.
\]

\section{物理理解}

有限脉冲的作用可以理解为：原子不再在三个无限短时刻 ``采样'' 激光相位，而是在每个 Bragg pulse 持续期间对激光相位进行带权平均。

因此：

\begin{itemize}[leftmargin=2em]
\item 当扰动频率满足
\[
\omega\ll\Omega_{\rm R},
\]
有限脉冲结果接近瞬时脉冲结果；
\item 当
\[
\omega\sim\Omega_{\rm R},
\]
有限脉冲效应开始明显；
\item Rabi 频率越高，脉冲越短，原子干涉仪越接近理想瞬时采样。
\end{itemize}

因此，有限 Rabi frequency 实际上给原子干涉仪引入了一个额外的高频响应尺度。

\section{最终计算公式总结}

给定
\[
f_{\rm R},\quad T,\quad \lambda,
\]
按以下顺序计算：

\[
\boxed{
\Omega_{\rm R}=2\pi f_{\rm R}
}
\]

\[
\boxed{
\tau=\frac{1}{4f_{\rm R}}
}
\]

\[
\boxed{
s(t)
\quad\Longrightarrow\quad
G(\omega)=\int s(t)e^{-i\omega t}dt
}
\]

\[
\boxed{
H_\varphi(\omega)=i\omega G(\omega)
}
\]

\[
\boxed{
|H_\varphi(\omega)|
=
|\omega G(\omega)|
}
\]

一阶 Bragg：
\[
\boxed{
k_{\rm eff}=2k=\frac{4\pi}{\lambda}
}
\]

位移输入：
\[
\boxed{
H_{x\rightarrow\Phi}
=
k_{\rm eff}i\omega G(\omega)
}
\]

加速度输入：
\[
\boxed{
H_{a\rightarrow\Phi}
=
-k_{\rm eff}
\frac{i\omega G(\omega)}{\omega^2}
}
\]

瞬时脉冲极限校验：
\[
\boxed{
|H_\varphi|
\rightarrow
4\sin^2\left(\frac{\omega T}{2}\right)
}
\]

以及
\[
\boxed{
|H_\varphi|^2
\rightarrow
16\sin^4\left(\frac{\omega T}{2}\right)
}.
\]

\section*{备注}

本文只讨论无光学谐振腔、两束反向传播平面波、一阶 Bragg、共振、方波脉冲以及两能级近似。若进一步考虑失谐、速度分布、多动量态、高阶 Bragg、Gaussian pulse、AC Stark shift 或 cavity pulse filtering，则需要从完整的 Bragg Hamiltonian 数值求解原子态演化，并据此重新计算 sensitivity function。

\end{document}
